\thispagestyle{empty}\begin{center}\vspace*{-.65in}{\Large\bfseries Carl Friedrich Gauss\par}\bigskip{\Large Works, Volume IV\par}\bigskip{\large Theory of the Combination of Observations Subject to the Smallest Errors\par}\medskip{\large First Part, Second Part and Supplement\par}\end{center}\smallskip\begingroup\small\setlength{\baselineskip}{11.0pt}\noindent AI-translated into English from the historical Latin and German. The German volume titles are AI-translated into English. This reader contains the historical volume titles, the formal First Part title and the complete First Part, Articles 1--22, including both author footnotes, and the formal and repeated titles of the complete Second Part, Articles 23--40, including the author footnote. The complete Supplement follows, with its formal and repeated titles, Articles 1--25 and all four author footnotes. The complete notices of 26 February 1821, 24 February 1823 and 25 September 1826 from the Göttingische gelehrte Anzeigen are also included, followed by the complete Determination of the Accuracy of Observations of 1816, Articles 1--8, with its handwritten note and numerical tables. The complete posthumous Section I, \textit{General overview of the method}, on the balance sheet of widows\textquotesingle\ funds, signed 9 January 1845 and including its author footnote, follows. The complete Section II, \textit{Investigation of the present condition of the professors\textquotesingle\ widows\textquotesingle\ fund at Göttingen}, follows, with its preface, both divisions, all author notes and the closing balance signed 2 November 1845. The complete Section III, the commission report with all four author footnotes and its accompanying note, follows. The complete balance report signed 19 October 1851 follows, including the nine balance items and the closing recommendations. The accompanying editorial introduction and all 134 entries of the membership table follow. The complete annuity work follows, including its title, all ten tables, explanatory notes, worked examples and Schering's closing remarks. The complete mapping paper follows, with both titles, its preface, Articles 1--14, the author footnote and the marginal remarks. The complete 1827 paper on curved surfaces follows, with both titles, all 29 articles and its author footnote. The complete first higher-geodesy memoir follows, with both titles, the preface, all 17 articles, both author footnotes, and the appended tables. The complete second higher-geodesy memoir follows, with both titles, the introduction, Articles 18--35 and the appended table. Complete notices cover curved surfaces (1827), higher geodesy (1830, 1843, 1846), Mollweide on Archimedes (1808), Monge and Hachette (1813), Herschel, Mollweide and Kries (1814), parallel lines (1816, 1822), and the French military journal (1830). Also included are the 1810 ellipse paper, the Carnot additions from \textit{Werke} and the geometrical solutions of 1823 and 1850, and the complete general register of coordinates with Schering’s remarks. The survey outlines (pp.~449--480), including the Danish tables, are followed by the degree-survey reports, plans and memoranda, the errata, contents and printer’s imprint. The contents retain the original page references.\par\medskip\noindent Historical author: Carl Friedrich Gauss. Volume issued by the Royal Society of Sciences at Göttingen, 1873. The First Part was presented on 15 February 1821 and published in 1823. Historical source text and portrait are public domain; translation and typesetting are AI-produced. The historical original is retained separately.\par\endgroup\clearpage
